
Benchmark contamination—when test data appears in training sets—invalidates model evaluation. Fu et al. (2024) demonstrated that membership inference attacks fail during pretraining, achieving AUC≈50\% because models learn distributions rather than memorizing instances. Dekoninck et al. (2024) showed that adversarial paraphrasing (the "EAL" attack) achieves ~15\% benchmark gains while evading semantic detection (TPR<2\% at 1\% FPR). These findings motivated investigation of geometric detection methods analyzing training trajectory signatures.

We designed a three-tier architecture: (1) data-layer filters using locality-sensitive hashing, (2) Task Signature Graph probes measuring differential alignment, and (3) geometric trajectory analysis tracking gradient alignment, Hessian anisotropy, CKA compression, and information efficiency. The hypothesis was that contamination-induced gains require detectable parameter-space alignment.

The experiment failed completely. The detector achieved 100\% detection power and 100\% false positive rate—the signature of a constant classifier. All three tiers returned hard-coded \texttt{True} values without computing metrics. Root cause: 10 of 15 implementation tasks (67\%) were not executed, targeting all detection components. This occurred despite passing validation that verified real datasets (GSM8K, MATH) were loaded correctly.

We report this failure for two reasons: (1) to document a validation methodology gap observed in this research system, and (2) to clarify that the original theoretical questions remain unanswered. The hypothesis that geometric signatures correlate with contamination was never tested.
